\section{Conclusion}
\label{sec:conclusion}
Our work reveals that the emergent phenomena of massive and artifact tokens in vision transformers govern the information flow through attention mechanisms and present an opportunity for efficiency gains. By introducing Fast Nystr\"om Attention (FNA), a training-free approach that exploits these token properties for linear-time, low-rank approximations of self-attention, we demonstrate significant reductions in computational and memory overhead while preserving competitive performance on a variety of downstream tasks. Our comprehensive analysis—spanning iterative and non-iterative detection methods as well as the strategic masking of attention sinks—sheds light on the underlying suppression dynamics that shape token interactions in attention, enabling us to enhance global feature aggregation and improve downstream tasks.

% While our analysis reveals that the behavior of massive and artifact tokens varies across different vision transformer architectures (and is even absent in models such as MAE and DINOv1) [reference Appendix], our findings lay a robust groundwork for understanding these mechanisms at inference. Determining the precise conditions under which these behaviors emerge during training remains an exciting avenue for future research, potentially leading to further enhancements in transformer performance and efficiency.


